\cite{mukherjee2025detailed} realizam uma análise comparativa detalhada entre as métricas de avaliação automática BLEURT e BERTScore, investigando suas forças e limitações em cenários específicos de tradução automática. O estudo utiliza uma metodologia baseada em três condições ideais para métricas de tradução automática, avaliando o comportamento das métricas em diferentes contextos semânticos. Os experimentos são conduzidos em três conjuntos de dados: TQ-AutoTest (com sentenças de aproximadamente 6 palavras, focado em fenômenos linguísticos), PE²rr (com sentenças de cerca de 22-23 palavras, contendo anotações de erros em nível de palavra) e GEC CoNLL (focado em correção de erros gramaticais). Além disso, o estudo discute a diferença entre métricas baseadas em \textit{embedding} BLEURT/BERTScore e métricas baseadas em sintaxe, como o \textit{BLEU}.

\cite{di2025meta} apresentam o corpus MT@BZ, um conjunto de dados de decretos traduzidos por máquina para os pares linguísticos Italiano-Alemão Sul-Tirolês, anotado manualmente com uma organização personalizada de erros de tradução. Sobre esse corpus,os autores realizaram uma análise qualitativa detalhada dos erros cometidos por sistemas de tradução automática no domínio jurídico. No estudo, os autores conduzem uma avaliação sistemática de métricas automáticas de avaliação de tradução (AMTE), investigando o desempenho de diferentes paradigmas, incluindo métricas baseadas em string (BLEU, chrF, TER), baseadas em embeddings (BERTScore com diferentes backbones) e métricas neurais aprendidas (COMET, BLEURT, UNITE, MetricX, XCOMET) na tarefa de avaliar a qualidade de traduções entre italiano e alemão sul-tirolês. Para estabelecer um padrão ouro de qualidade, os autores atribuíram pesos de classificação a cada tipo de erro anotado no corpus, com base em consulta a especialistas em direito e linguística do contexto sul-tirolês, permitindo a comparação entre as pontuações automáticas e os julgamentos humanos.

\cite{xu2309paradigm} apresenta um estudo de caso sobre a criação de recursos de tradução para cinco línguas de Uganda(Luganda, Runyankole, Acholi, Lugbara e Ateso) utilizando o modelo criado a partir do framework \textit{MarianMT} através de técnicas de refinamento(\textit{fine-tuning}).

\cite{moslem2026afrinllb} apresenta a proposta de AfriNLLB, uma série de modelos leves e eficientes para tradução automática envolvendo línguas africanas construídos a partir do modelo pré treinado \textit{NLLB} utilizando bases de dados públicas para geração de novos recursos na área de tradução automática.

Dado os trabalhos apresentados, eles serão utilizados para montar a base metodologica deste trabalho, definindo elementos, tais como: modelos, tipos de dados e métricas de avaliação que serão utilizados ao decorrer do trabalho.

